\begingroup
\tiny
\renewcommand{\arraystretch}{1.05}
\setlength{\LTleft}{-1.0cm}
\setlength{\LTright}{0pt}

\begin{longtable}{p{0.9cm} p{6cm} p{0.9cm} p{0.9cm} p{6cm}}
\caption{Codierte Auswertungsmatrix der Experteninterviews}
\label{tab:codierte-auswertungsmatrix}\\

\toprule
\textbf{Fundstelle} & \textbf{Paraphrase der Aussage} & \textbf{Primär\-code} & \textbf{Sekun\-där\-code} & \textbf{Analytische Notiz / Modellfolge} \\
\midrule
\endfirsthead

\toprule
\textbf{Fundstelle} & \textbf{Paraphrase der Aussage} & \textbf{Primär\-code} & \textbf{Sekun\-där\-code} & \textbf{Analytische Notiz / Modellfolge} \\
\midrule
\endhead

\midrule
\multicolumn{5}{r}{\textit{Fortsetzung auf der nächsten Seite}} \\
\midrule
\endfoot

\bottomrule
\endlastfoot

\multicolumn{5}{l}{\textbf{Interview I1}} \\
\midrule
I1\_002 & DCAT-AP.de-Validität wird als Mindestvoraussetzung hochwertiger GovData-Metadaten beschrieben. & K3.1 & K1.1 & Basiskonformität als Mindestanforderung im Modell ausweisen. \\
I1\_003 & Best-Practice-Empfehlungen und Beispieldatensätze markieren eine höhere Qualitätsstufe jenseits der Mindestvalidität. & K1.1 & K2.2 & Pflicht- und Best-Practice-Erfüllung getrennt bewerten. \\
I1\_004 & Zusätzliche Detailtiefe steht in einem Aufwand-Nutzen-Verhältnis mit abnehmendem Grenznutzen. & K1.2 & K7.5 & Score nicht durch maximale Feldanzahl überladen. \\
I1\_006 & Datenqualität wird als kontextsensitiv und zweckabhängig beschrieben; GovData soll Auffindbarkeit und Zugang erleichtern. & K7.1 & K1.3 & Bewertung an Nutzungskontext und Such-/Zugriffsfunktion ausrichten. \\
I1\_007 & Reichhaltigkeit und Vollständigkeit der Metadaten gelten als zentrale Qualitätskriterien. & K2.2 & K2.1 & Empfohlene Felder und Reichhaltigkeit im Score berücksichtigen. \\
I1\_009 & Optionale Rollen wie Creator oder Maintainer sind je nach Domäne unterschiedlich relevant. & K7.1 & K2.3 & Domänenspezifische Feldgewichtung ermöglichen. \\
I1\_010 & Allgemeine DCAT-AP.de-Eigenschaften sind breit relevant, spezifische Relationen stärker kontextabhängig. & K7.1 & K2.4 & Unterscheidung zwischen allgemeinen und kontextspezifischen Indikatoren. \\
I1\_012 & Föderale Portale unterscheiden sich in Leistungsfähigkeit und Veröffentlichungsprozessen. & K7.3 & K9.4 & Vergleichbarkeit zwischen Portalen vorsichtig kontextualisieren. \\
I1\_013 & Kontextsensitive Beispieldatensätze werden genutzt, um thematische Vielfalt sichtbar zu machen. & K7.1 & K4.3 & Evaluation mit thematisch unterschiedlichen Datensätzen planen. \\
I1\_014 & Schleswig-Holstein wird wegen vorgeschalteter Qualitätssicherung und Best-Practice-Umsetzung als positiver Referenzpunkt genannt. & K1.1 & K9.4 & Gute Beispiele als Benchmark für Modelltests nutzen. \\
I1\_015 & Schlecht gepflegte, veraltete, unklare oder nicht erreichbare Datensätze werden als typische Qualitätsprobleme beschrieben. & K5.1 & K4.4 & Verfügbarkeit und Beschreibungsklarheit als zentrale Prüffelder einbauen. \\
I1\_017 & Bestehende Scores sollten stärker an etablierten Datenqualitätskriterien und Fit-for-Purpose-Konzepten anknüpfen. & K9.1 & K7.1 & Indikatorenauswahl theoretisch begründen. \\
I1\_018 & Klassische Datenqualitätskriterien werden als leicht nutzbare Grundlage für Analysewerkzeuge beschrieben. & K9.1 & K3.5 & Literaturbasierte Basisdimensionen in das Modell integrieren. \\
I1\_019 & Funktionierende Links, Fehlerfreiheit, Alter, Aktualisierung und Harvesting-Zeitpunkt sind relevante Basiskriterien. & K5.3 & K7.5 & Aktualität und Harvesting als eigene Kontextindikatoren prüfen. \\
I1\_020 & Linked-Data-Qualität wird über intensive RDF-Nutzung, Vokabularverweise und IDs statt Literale beschrieben. & K6.1 & K6.3 & Linked-Data-Reife als Qualitätsdimension aufnehmen. \\
I1\_022 & Eine GND-ID wird als verlässlicher bewertet als derselbe Name als Freitextliteral. & K6.1 & K6.2 & Normdatenbezug als Verlässlichkeitsindikator operationalisieren. \\
I1\_024 & Harvesting- und Katalog-Metadaten werden als interne Qualitätskriterien für GovData beschrieben. & K5.3 & K9.1 & Catalog-Record-Informationen als mögliche Zusatzdimension prüfen. \\
I1\_025 & Catalog Records könnten Validierungen, Harvesting-Herkunft und Aktualisierungsinformationen dokumentieren. & K5.3 & K6.3 & Katalog-Metadaten für Qualitätsnachweise nutzen. \\
I1\_027 & DCAT-AP.de-Validierung wird als wichtige Basis unter Fehlerfreiheit eingeordnet. & K3.2 & K9.2 & SHACL/Validierung als formale Basiskomponente verwenden. \\
I1\_028 & Validierung ist wichtig, aber nicht ausreichend; Syntax, formale und inhaltliche Qualität benötigen unterschiedliche Methoden. & K9.2 & K4.4 & Formale und semantische Teildimensionen getrennt ausweisen. \\
I1\_029 & DCAT-AP.de-Qualität soll als Dimension im Gesamtscoring angezeigt werden, Gewichtung bleibt schwierig. & K9.2 & K3.2 & SHACL nicht als Gesamtscore interpretieren. \\
I1\_031 & Stark normierte Ausgangsdaten, insbesondere Geodaten, erleichtern hochwertige Metadaten. & K7.2 & K8.5 & Domänenspezifische Ausgangslage bei Bewertung berücksichtigen. \\
I1\_032 & Die Grenze zwischen Daten und Metadaten ist fließend; normierte Bereiche bieten reichhaltige Informationen. & K8.5 & K7.2 & Rohdatenbezug domänenspezifisch begrenzen. \\
I1\_033 & Freie, einzigartige oder manuell erhobene Daten sind schwieriger gut zu beschreiben. & K7.1 & K4.4 & Beschreibungsgüte abhängig vom Datentyp interpretieren. \\
I1\_035 & Freitext und natürliche Sprache bieten Potenzial für LLM-/KI-gestützte Fehler- und Schwächenanalyse. & K10.2 & K3.5 & KI für Freitext- und einfache Plausibilitätsprüfungen prüfen. \\
I1\_036 & KI kann fehlertolerant unterstützen, ersetzt aber teilweise nur klassisch mögliche Verfahren. & K10.1 & K10.9 & KI-Einsatz nur bei zusätzlichem Nutzen rechtfertigen. \\
I1\_037 & Kontextsensitive semantische KI-Analyse kann fachliche Begriffe fehlinterpretieren und Schaden erzeugen. & K10.8 & K4.3 & Keine autonome Korrektur fachsprachlicher Metadaten. \\
I1\_039 & Eine agentische KI könnte als Metadaten-Linter auf Basis von DCAT-AP.de, SHACL und Soll-/Kann-Kriterien dienen. & K10.1 & K3.2 & Linter-Funktion als unterstützendes Analysemodul konzipieren. \\
I1\_040 & KI kann Lücken identifizieren und föderale Vergleichbarkeit zwischen ähnlichen Datensätzen unterstützen. & K7.3 & K10.6 & Vergleichsanalyse als KI-unterstützte Kontextfunktion. \\
I1\_041 & Assisted Learning kann Freitext-Stichworte auf Kategorien, Schlagworte oder EU-Vokabulare mappen. & K10.3 & K6.3 & Keyword-Harmonisierung als potenzieller KI-Indikator. \\
I1\_043 & KI ist bei natürlicher Sprache, einfacher Sprache, Barrierearmut, Fachtermini und Synonymen sinnvoll. & K10.2 & K4.3 & Verständlichkeits- und Synonymprüfung als KI-gestützte Ergänzung. \\
I1\_044 & KI kann Beschreibungstexte anhand von Style Guides prüfen und Verbesserungsvorschläge machen. & K10.2 & K9.5 & KI nur als Vorschlags- und Verbesserungsebene nutzen. \\
I1\_046 & Generative KI ist für Scoring problematisch, da stochastische Ausreißer und hoher Wiederholungsaufwand entstehen. & K10.8 & K10.9 & KI nicht als autonome Bewertungsinstanz einsetzen. \\
I1\_047 & Qualitätsfragen erfordern häufig spezifische SPARQL-Abfragen; Stakeholder haben unterschiedliche Bedarfe. & K10.4 & K7.4 & SPARQL-Assistenz für domänenspezifische Qualitätsfragen prüfen. \\
I1\_048 & Regionale Datentypen benötigen angepasste Queries und spezifische Analysebausteine. & K10.4 & K7.1 & Templates oder KI-Assistenz für kontextspezifische Abfragen nutzen. \\
I1\_049 & KI kann Fachpersonen ohne tiefe technische Expertise Zugang zu Auswertungen geben. & K10.4 & K7.4 & Assistenzsystem statt autonomem Score-Generator vorsehen. \\
\multicolumn{5}{l}{\textbf{Interview I2}} \\
\midrule
I2\_002 & Nicht perfekte Metadaten werden als häufiges Praxisphänomen beschrieben. & K1.1 & K9.5 & Modell muss typische Fehler verständlich erfassen. \\
I2\_003 & Perfekte Metadaten basieren idealerweise auf DCAT-AP.de beziehungsweise einem semantischen Graphen. & K1.4 & K6.4 & Graphfähigkeit als Idealstruktur berücksichtigen. \\
I2\_004 & Pflichtfelder sind notwendig; fachlich vorhandene Informationen sollten zusätzlich übernommen werden. & K2.3 & K2.2 & Feldreichhaltigkeit über Pflichtfelder hinaus bewerten. \\
I2\_005 & Semantische Modellierung, externe Referenzen, Wikidata, GND und Thesauri werden als Ideal genannt. & K6.2 & K6.3 & Externe Referenzierung und Thesaurusnutzung prüfen. \\
I2\_006 & Querverbindungen zwischen vergleichbaren kommunalen Datensätzen werden als Idealfall beschrieben. & K6.4 & K7.3 & Interlinking zwischen vergleichbaren Datensätzen als Qualitätsaspekt. \\
I2\_008 & Zu kurze Titel und identische Titel/Beschreibungen machen Metadaten unverständlich. & K4.1 & K4.2 & Titelqualität und Redundanzprüfung aufnehmen. \\
I2\_009 & Regionale Begriffe wie „I-Dötzchen“ erschweren die Verständlichkeit. & K4.3 & K10.2 & Regionale Begriffe und Synonyme als Freitextproblem berücksichtigen. \\
I2\_010 & Fachsprache kann detailliert und korrekt, aber für Außenstehende unverständlich sein. & K4.3 & K4.4 & Verständlichkeitsprüfung nicht mit fachlicher Korrektheit verwechseln. \\
I2\_011 & Dead Links, zu wenige Metadaten und widersprüchliche Angaben werden als schlechte Metadaten beschrieben. & K5.1 & K4.5 & Verfügbarkeit und Plausibilität getrennt prüfen. \\
I2\_012 & Eigentliche Daten sind oft aussagekräftiger als menschlich eingetragene Metadaten. & K8.3 & K8.5 & Rohdaten als ergänzende Prüfbasis berücksichtigen. \\
I2\_013 & Bounding Boxes sind manuell eingetragen und fehleranfällig; Berechnung aus Daten wäre sinnvoller. & K8.4 & K7.2 & Geobezogene Angaben aus Daten ableiten oder validieren. \\
I2\_014 & Metadatenfehler entstehen häufig, weil Bereitstellende Zweck und Nachnutzung nicht verstehen. & K1.3 & K9.5 & Reporting muss Zweck und Nutzen von Feldern erklären. \\
I2\_016 & Praxiskenntnis vieler guter und schlechter Datensätze wird bestätigt. & K9.4 & K1.1 & Evaluation mit Beispieldatensätzen stützen. \\
I2\_017 & Portalübergreifende Verarbeitung kann Dubletten und technische Verschlechterungen erzeugen. & K6.4 & K5.2 & Provenienz und Datensatzidentität stärker berücksichtigen. \\
I2\_020 & Diskussionen zu Scores betreffen besonders wissenschaftliche und semantische Grundlagen. & K9.1 & K4.6 & Kriterien theoretisch und semantisch begründen. \\
I2\_021 & Aktuelle Scores prüfen oft ohne ausreichend begründete wissenschaftliche Grundlage. & K9.1 & K1.1 & Kriterienauswahl im Modell transparent begründen. \\
I2\_022 & EU- und Schweizer Scores werden wegen fehlender Begründung trotz FAIR-Bezug kritisiert. & K9.1 & K2.2 & Übernommene Metriken nicht unreflektiert verwenden. \\
I2\_023 & GovData will Qualität sinnvoll messen; was gemessen wird, muss geklärt sein. & K9.1 & K9.3 & Indikatoren vor Scoringlogik konzeptionell begründen. \\
I2\_024 & Reine Anzahl von Datensätzen ist als Erfolgs- oder Qualitätsmetrik ungeeignet. & K9.1 & K1.2 & Keine quantitätsgetriebene Portallogik als Qualitätsersatz. \\
I2\_026 & Absolute Vollständigkeit ist im Graphen unmöglich; relevant ist relative Vollständigkeit je Fachthema. & K2.1 & K7.1 & Vollständigkeit relational und domänenspezifisch interpretieren. \\
I2\_027 & Sprachliche Verständlichkeit ist zentral; Länge allein reicht nicht. & K4.4 & K10.2 & Beschreibungsgüte semantisch statt nur formal prüfen. \\
I2\_028 & Distributionen sollten verschiedene Formate derselben Daten korrekt beschreiben. & K5.2 & K8.1 & Distributionen und Formatbeschreibung prüfen. \\
I2\_029 & Blank Nodes vermeiden, Querverweise nutzen und Organisationen strukturiert abbilden. & K6.5 & K6.4 & Strukturelle Linked-Data-Qualität einbeziehen. \\
I2\_030 & Attribution und korrekte Quellenangabe bei Lizenzen gehören zur Metadatenqualität. & K4.6 & K7.4 & Lizenz- und Attributionsangaben als Reuse-Kontext prüfen. \\
I2\_031 & Der Weg eines Datensatzes durch verschiedene Portale ist derzeit nicht sauber abbildbar. & K5.3 & K6.4 & Provenienz/Harvesting-Pfad als Zusatzinformation erwägen. \\
I2\_032 & Verbesserung von Metadaten durch Portale wirft Fragen der Datenhoheit auf. & K10.7 & K9.5 & Automatische Korrektur von Analyse und Empfehlung trennen. \\
I2\_034 & SHACL-/Qualitätsrankings können für Ländervertretungen und Bund nützlich sein. & K9.4 & K9.2 & Benchmarking als Verbesserungsinstrument möglich. \\
I2\_035 & Länder und Kommunen möchten sich vergleichen und Ressourcenbedarf sichtbar machen. & K9.4 & K7.3 & Vergleiche kontextualisiert bereitstellen. \\
I2\_036 & Rankings sollen nicht bloßstellen, sondern Hinweise auf Verbesserungsmöglichkeiten geben. & K9.5 & K9.4 & Reporting handlungsorientiert formulieren. \\
I2\_037 & Regelmäßige Prüfungen können dazu führen, dass Fehler behoben werden. & K9.4 & K5.3 & Wiederholte Prüfung als Monitoringfunktion einplanen. \\
I2\_039 & Teil-Scores helfen, Problembereiche zu lokalisieren. & K9.3 & K9.5 & Mehrdimensionale Teilscores statt nur Gesamtscore. \\
I2\_040 & SHACL-Ergebnisse sind häufig technisch und müssen interpretiert werden. & K9.5 & K3.2 & Validatorausgaben in verständliches Reporting übersetzen. \\
I2\_041 & Fehlermeldungen müssen erklären, was konkret zu tun ist. & K9.5 & K3.4 & Handlungshinweise pro Fehlerklasse entwickeln. \\
I2\_043 & Grenze zwischen Metadaten und Daten verschwimmt im GovData-Kontext. & K8.5 & K1.4 & Rohdatenbezug methodisch abgrenzen. \\
I2\_044 & Geodatenservices sind APIs, die sich selbst beschreiben und Daten unter der Metadatenebene enthalten. & K8.2 & K7.2 & Services anders bewerten als statische Dateien. \\
I2\_045 & Metadaten und Daten korrelieren; falsche Formatangaben verletzen Metadatenqualität. & K8.1 & K5.4 & Formatabgleich als automatisierbarer Indikator. \\
I2\_046 & Erwartetes und tatsächliches Serviceverhalten kann auseinanderfallen, etwa Karte versus XML-Fehlermeldung. & K8.2 & K5.2 & Service-Output gegen Nutzungserwartung prüfen. \\
I2\_047 & Titel und tatsächlicher Dateninhalt können nicht zusammenpassen. & K8.3 & K4.5 & Inhaltliche Metadaten-Daten-Kongruenz prüfen. \\
I2\_049 & AI Agents sollten nicht von vornherein in den Vordergrund gestellt werden. & K10.8 & K10.1 & KI-Einsatz fachlich begründen. \\
I2\_050 & Vor KI-Einsatz ist zu prüfen, ob andere Mittel dasselbe Problem lösen können. & K10.9 & K9.1 & KI nur bei methodischem Mehrwert einsetzen. \\
I2\_051 & AI Agents sind nicht automatisch intelligent; entscheidend ist die Art der Ansteuerung. & K10.8 & K10.1 & Automatisierungsgrad begrenzen. \\
I2\_052 & Textanalyse ist schwach; KI auf semantischem Graphen ist verlässlicher und halluziniert weniger. & K10.5 & K10.8 & Graphbasierte KI statt rein textbasierter KI bevorzugen. \\
I2\_054 & Agenten könnten Graphen, europäische Portale und wissenschaftliche Portale als Kontext heranziehen. & K10.5 & K6.2 & Externe Kontexte kontrolliert einbeziehen. \\
I2\_055 & Schlagwörter sind oft Zeichenketten; Thesauri können zusätzlichen Kontext liefern. & K6.3 & K10.3 & Keyword-Kontextualisierung über Vokabulare unterstützen. \\
I2\_056 & Verwandte Konzepte im Graphen verbessern Suche und Bewertung. & K6.4 & K10.5 & Graphbasierte semantische Kontextanalyse ermöglichen. \\
I2\_058 & Portale sollten stärker in die Daten selbst hineinschauen, um Passung zu prüfen. & K8.3 & K10.6 & Rohdatenprüfung als Erweiterung des Scorings. \\
I2\_059 & Agenten könnten Strukturen, Vergleichbarkeit und Flächendeckung prüfen. & K10.6 & K8.4 & KI für komplexe Rohdatenmuster prüfen. \\
I2\_060 & KI kann große Datenmengen, ältere Sprachmuster und historische Kontexte analysieren. & K10.6 & K4.3 & KI für Massenauswertung und Sprachkontext nutzbar. \\
I2\_061 & Häufigkeiten in großen Datensatzmengen sind manuell kaum analysierbar. & K10.6 & K9.3 & Skalierende Analysebausteine vorsehen. \\
I2\_062 & KI erzeugt Serverlast und benötigt Training sowie Betrieb. & K10.9 & K10.8 & Betriebskosten und Infrastruktur als Grenze einplanen. \\
I2\_063 & KI ist sinnvoll, aber nicht alles muss mit KI gelöst werden. & K10.9 & K9.1 & Aufwand-Nutzen-Prüfung in KI-Konzept aufnehmen. \\
\multicolumn{5}{l}{\textbf{Interview I3}} \\
\midrule
I3\_002 & Perfekte Metadaten enthalten relevante Felder in hoher inhaltlicher und syntaktischer Qualität; nicht alle Felder sind immer sinnvoll. & K1.2 & K2.4 & Relevanzbasierte statt maximale Feldbefüllung. \\
I3\_003 & Wichtige Felder sollen sorgfältig, syntaktisch korrekt und nachnutzbar ausgefüllt sein. & K2.4 & K3.3 & Syntaktische Korrektheit und Nachnutzbarkeit verbinden. \\
I3\_005 & Metadaten sollen Inhalts- und technische Vorschau bieten, bevor der Datensatz geöffnet wird. & K1.3 & K4.4 & Nutzungsentscheidung durch Metadaten ermöglichen. \\
I3\_006 & Konsistenz und Standardisierung sind bei ähnlichen Daten wichtig, etwa Einwohnerstatistiken. & K4.6 & K6.3 & Synonym- und Standardisierungsprobleme berücksichtigen. \\
I3\_008 & Schlechte Metadaten finden sich sowohl auf GovData als auch im eigenen Portal. & K1.1 & K9.4 & Breite Testbasis für Prototyp möglich. \\
I3\_010 & piveau prüft Feldexistenz, Validator Syntax; inhaltliche Sinnhaftigkeit bleibt unberücksichtigt. & K4.4 & K9.1 & Semantische Qualität als Ergänzung bestehender Ansätze. \\
I3\_012 & Scoring sollte mehrstufig prüfen: Existenz, syntaktische Korrektheit, semantische und inhaltliche Passung. & K9.3 & K3.4 & Mehrstufige Bewertungsarchitektur ableiten. \\
I3\_013 & Pflichtfelder allein reichen nicht; auch über DCAT-Pflichtfelder hinausgehende Felder sind relevant. & K2.2 & K2.4 & Empfohlene und weitere relevante Felder berücksichtigen. \\
I3\_015 & Validierung verbindet Feldexistenz mit syntaktisch korrekter Ausfüllung. & K3.2 & K3.4 & SHACL als Verbindung von Vollständigkeit und Syntax nutzen. \\
I3\_017 & Validierungsergebnisse sollen klar trennen, ob ein Feld fehlt oder syntaktisch fehlerhaft ist. & K3.4 & K9.5 & Fehlerarten explizit im Report ausweisen. \\
I3\_018 & Valides XML, korrekte URIs und Vokabulare sind grundlegende Prüfungen. & K3.5 & K3.3 & Basale technische Prüfungen als harte Kriterien. \\
I3\_020 & Rohdaten sind für Gegenprüfung relevant; falsche Formatangaben zeigen Metadatenfehler. & K8.1 & K5.4 & Formatkongruenz als automatisierbarer Indikator. \\
I3\_021 & Rohdatenqualität betrifft nicht-proprietäre und gut lesbare Formate, wird aber auch an Bereitstellende herangetragen. & K5.4 & K8.5 & Rohdatenqualität nur angrenzend berücksichtigen. \\
I3\_022 & Metadatenportale können Rohdaten fachlich nur begrenzt bewerten; Fachstandards unterscheiden sich. & K8.5 & K7.2 & Fachliche Grenzen des Modells transparent machen. \\
I3\_024 & Kontextsensitive KI kann Rohdaten und Umfeld nutzen, um Metadaten über Syntax hinaus zu prüfen. & K10.6 & K8.3 & KI-gestützte Kontextprüfung als Erweiterung prüfen. \\
I3\_026 & Große Mengen könnten schneller bewertet werden; Datenhoheit bei Veränderung bleibt kritisch. & K10.7 & K10.9 & Messung von automatischer Veränderung trennen. \\
I3\_028 & Format, Dateigröße, zeitliche Bezüge, Titel, Beschreibung, Keywords, Themenfeld und Kategorie sind relevante Prüfinformationen. & K8.1 & K4.6 & Technische und semantische Kontextfelder operationalisieren. \\
I3\_029 & Lizenz, Publisher und Contact Point sind wichtig, aber schwer kontextsensitiv automatisiert zu prüfen. & K8.5 & K7.4 & Rechtliche und Herkunftsfelder vorsichtig bewerten. \\
I3\_031 & Weitere KI-Anwendungsfälle werden nicht konkret benannt. & K10.9 &  & Kein zusätzlicher Indikator; als negative Evidenz dokumentieren. \\

\end{longtable}